
Projected on the numerical distribution function grid, the electron dynamics
is expressed 
\begin{eqnarray}
&&\frac{\widetilde{q}\left( \psi _{l+1/2}\right) }{B_{\min }\left( \psi
_{l+1/2}\right) }p_{i+1/2}^{2}\frac{f_{0,l+1/2,i+1/2,j+1/2}^{\left( 0\right)
(k+1)}-f_{0,l+1/2,i+1/2,j+1/2}^{\left( 0\right) (k)}}{\Delta t}  \nonumber \\
&=&\frac{\widetilde{q}\left( \psi _{l+1/2}\right) }{B_{\min }\left( \psi
_{l+1/2}\right) }\sum_{i^{\prime }=i-1}^{i^{\prime }=i+1}\sum_{j^{\prime
}=j-1}^{j^{\prime }=j+1}\overline{\overline{M}}_{p,l+1/2,i^{\prime
}+1/2,j^{\prime }+1/2}^{\left( 0\right) }f_{0,l+1/2,i^{\prime
}+1/2,j^{\prime }+1/2}^{\left( 0\right) (k+1)}  \nonumber \\
&&+\frac{p_{i+1/2}^{2}}{\lambda \left( \psi _{l+1/2},\xi _{0,j+1/2}\right) }%
\sum_{l^{\prime }=l-1}^{l^{\prime }=l+1}\overline{\overline{M}}_{\psi
,l^{\prime }+1/2,i+1/2,j+1/2}^{\left( 0\right) }f_{0,l^{\prime
}+1/2,i+1/2,j+1/2}^{\left( 0\right) (k+1)}
\end{eqnarray}
where the first right handside term represents the dynamics in momentum
space while the second one concerns the spatial dynamics. Here upwind
differencing is performed, that corresponds to the fully implicit scheme
that is usually numerically stable for arbitrary values of $\Delta t.$ 

In matrix form, 
\[
I
\]

By definition, $\overline{\overline{M}}_{p}^{\left( 0\right) }$ corresponds
to 
